\documentclass[11pt,a4paper]{article}

\usepackage[T1]{fontenc}
\usepackage[utf8]{inputenc}

\setlength{\parindent}{0pt}
\setlength{\parskip}{0.35em}

\newcommand{\placeholder}[1]{\textit{[#1]}}
\newenvironment{compactenum}
    {\begin{enumerate}\setlength{\itemsep}{0pt}\setlength{\parsep}{0pt}\setlength{\topsep}{0.25em}}
    {\end{enumerate}}

\begin{document}

\begin{center}
    {\Large\bfseries Assessment Criteria for TSRT08}
\end{center}

\section{Intended Learning Outcomes}

Optimal control consists of theory and methods for controlling dynamical systems so that a mathematical criterion is minimized. The course aims at giving basic theory and engineering oriented computational methods in the area. After taking the course the participants are expected to be able to:

\begin{compactenum}
    \item Analyze and synthesize optimal open loop control signals using the Maximum Principle.
    \item Analyze and synthesize optimal feedback laws using Dynamic Programming.
    \item Analyze and synthesize optimal feedback laws using Model Predictive Control.
    \item Analyze and synthesize optimal feedback laws using Reinforcement Learning.
    \item Use numerical software to solve optimal control problems.
\end{compactenum}

\section{Assessment Criteria}

For Grade 3, with respect to the first two learning outcomes, students should understand the following sections in the course book:

\begin{compactenum}
    \item 5.5, 8.1, 8.3, 8.4.
    \item 7.2, 7.3.
\end{compactenum}

For Grade 5, students should also understand the following sections in the course book:

\begin{compactenum}
    \item 8.2, 8.5, 8.10.1.
    \item 7.1, 7.4.
\end{compactenum}

Learning Outcomes 3--5 are assessed in the homework assignments. For Grade 4, all Grade 3 assessment criteria must be fulfilled, as well as most of the Grade 5 assessment criteria.

\newpage

\section{Operational Grading Criteria for TEN1\\Written Exam}

The exam consists of four questions, each worth 10 points. Questions 1 and 2 examine Learning Outcome 1, while Questions 3 and 4 examine Learning Outcome 2.

To obtain Grade 3, a student must score at least 15 points in total, including at least 6 points from Questions 1 or 2 and at least 6 points from Questions 3 or 4. The total score required for Grade 4 is 23 points, and for Grade 5 it is 30 points.

If, during marking, either of the 6-point requirements is not met, the remaining exam questions will not be marked.

\end{document}
